\begin{helper}
Let \(C\) be a fixed history cell, let \(F\subseteq C\) be one fixed-count
event, and let \(E\) be the target event.  Suppose \(0\le p\le1\) and
\(B\ge0\).  Let \(\{H_b\}_{b\in\mathcal H}\) be a finite family of
first-color branch events such that each \(H_b\) is contained in \(C\), the
events \(H_b\) are pairwise disjoint, and every sample in \(C\cap F\) lies in
one of the branches \(H_b\).  If, for every branch \(b\),
\[
\Pr(E\cap F\cap H_b)\le B\,\Pr(H_b),
\]
where probabilities are the finite sample weights from
\(\noderef{TwoBiteEventProbability}\), then
\[
\noderef{TwoBiteConditionalProbability}(E\cap F\mid C)\le B.
\]
\end{helper}
\begin{proof}
Write \(w(\omega)=\noderef{TwoBiteSampleWeight}(\omega)\).  The assumptions
\(0\le p\le1\) make every red-edge factor, blue-edge factor, and injection
factor in \(w(\omega)\) nonnegative, so every event probability is a finite
sum of nonnegative terms.  In particular probability is monotone under
event inclusion.

If \(\Pr(C)=0\), then the definition of
\(\noderef{TwoBiteConditionalProbability}\) gives
\(\Pr(E\cap F\mid C)=0\), and the desired inequality follows from \(B\ge0\).
Assume from now on that \(\Pr(C)>0\).

Because each \(H_b\subseteq C\), the branch events are disjoint, and every
point of \(C\cap F\) lies in some \(H_b\), the events
\[
E\cap F\cap H_b\qquad (b\in\mathcal H)
\]
are pairwise disjoint and their union contains \(E\cap F\cap C\).  Hence
finite additivity and monotonicity give
\[
\Pr(E\cap F\cap C)
\le
\sum_{b\in\mathcal H}\Pr(E\cap F\cap H_b).
\]
Apply the branchwise hypothesis to each summand:
\[
\sum_{b\in\mathcal H}\Pr(E\cap F\cap H_b)
\le
B\sum_{b\in\mathcal H}\Pr(H_b).
\]
Since the \(H_b\) are disjoint subsets of \(C\), their union is a subset of
\(C\), so another use of finite additivity and monotonicity gives
\[
\sum_{b\in\mathcal H}\Pr(H_b)
=\Pr\!\left(\bigcup_b H_b\right)
\le \Pr(C).
\]
Multiplying by the nonnegative number \(B\) preserves the inequality, and
therefore
\[
\Pr(E\cap F\cap C)\le B\,\Pr(C).
\]
Dividing by the positive denominator \(\Pr(C)\) gives
\[
\Pr(E\cap F\mid C)
=\frac{\Pr(E\cap F\cap C)}{\Pr(C)}
\le B.
\]
Thus summing over first-color branches costs no factor \(|\mathcal H|\); the
branch probabilities themselves are the averaging weights.
\end{proof}
